\SourceLeaf{491}%
Skagen, it could undermine the confidence that the results should command if those observations had been obtained with instruments that differed not only as individual instruments but even in kind. It will therefore be best to make the observations here with the meridian circle and also with a zenith sector, specifically with the \textit{very same} zenith sector used to observe at all the other places.

The zenith sector used for the Danish determinations at Lauenburg, Lysabbel and Skagen, and also at Copenhagen, in 1819 and 1820 is the same instrument with which the late General \textsc{Mudge} observed during the English degree survey in 1802. English and French scholars jointly made a series of observations with the very same instrument in 1818 at Dunkirk, the northernmost endpoint of the French degree survey. The English government then lent it for the Danish degree survey for 2 years. All these circumstances make it all the more desirable that this instrument should also be used in the Hanoverian degree survey. It will be especially important and instructive to compare it with our Reichenbach meridian circle by means of observations made at the same time and in one place.

There can be no doubt that this instrument will be lent for the Hanoverian degree survey with the same readiness as for the Danish one. — — — I add the following circumstances. This zenith sector, made by \textsc{Ramsdon}, belongs to the Board of Ordnance and is ultimately under the Duke of \textsc{Wellington}, and, following the recent death of General \textsc{Mudge}, indirectly under Lieutenant Colonel \textsc{Colby}.

It follows from the foregoing that the time when the larger operations involved in the degree survey that require special arrangements will begin cannot at present be stated with certainty. In any event, nothing further can be done in the course of the present year beyond some of the astronomical observations at the observatory here (which I have already begun) and at most some preliminary work in the surrounding area.

Göttingen, August 11, 1820.
\begin{flushright}C. F. \textsc{Gauss}.\end{flushright}

\par\bigskip
\pdfbookmark[0]{Memorandum on Court Councillor Gauss’s presence at part of the Danish degree survey.}{pm-gradmessung}
\begin{center}\large MEMORANDUM ON THE PRESENCE OF COURT COUNCILLOR GAUSS AT PART OF THE OPERATIONS OF THE DANISH DEGREE SURVEY.\end{center}
\par
I have previously had the honor of submitting a most respectful report to Your — — on the degree survey ordered in the royal Danish territory and entrusted to the Danish astronomer \textsc{Schumacher}, on its purpose, on the benefits that science is entitled to expect from it, and on the means employed in it. It gives me all the more pleasure that another astronomer, Dr. \textsc{Olbers}, who also witnessed the operations at Lauenburg last year, has just published views of this undertaking that agree entirely with mine in the journal Allgemeine Geographische Ephemeriden (vol. 17, issue 3). My own judgment of the Danish astronomer as a person might perhaps appear biased inasmuch as he previously (1809) completed his training at the university here.
